%% Hand-authored circuitikz figure -- NOT generated by maestro2tex.
%% GENERIC textbook charge-redistribution SAR architecture. This is deliberately
%% NOT a transcription of the implemented cell: no netlist or OA view was read to
%% draw it, and no device dimensions appear anywhere (they are proprietary).
%% Only the external port names match the block.
\begin{figure}[htbp]
  \centering
  \providecommand{\sarbv}[1]{\ensuremath{\langle #1\rangle}}
  \resizebox{\linewidth}{!}{%
\ctikzset{bipoles/length=0.9cm, capacitors/scale=0.85, resistors/scale=0.85, amplifiers/scale=1.35}
\begin{circuitikz}[
    every node/.append style={font=\footnotesize},
    nlab/.style={font=\scriptsize, inner sep=1.5pt},
    port/.style={font=\scriptsize\ttfamily, inner sep=2pt},
    note/.style={font=\scriptsize, align=left, inner sep=2pt},
    railnode/.style={circle, fill=black, inner sep=1.05pt},
    bus/.style={line width=1.4pt},
  ]

% ---- geometry -------------------------------------------------------
\def\ytop{0}     \def\ybot{-2.3}  \def\ypiv{-2.95}
\def\ycon{-3.75} \def\yref{-4.75}
\def\xsl{-4.3}   \def\xsr{-3.0}   \def\xrr{10.9}
\def\xB{0} \def\xC{1.95} \def\xD{3.9} \def\xE{5.85} \def\xF{7.8} \def\xG{9.75}

% ===== sampling switch and top-plate rail ============================
\draw (\xsl-0.8,\ytop) -- (\xsl,\ytop);
\node[port, anchor=east] at (\xsl-0.9,\ytop){Analog\_Input};
\draw (\xsl,\ytop) to[nos] (\xsr,\ytop);
\node[nlab, anchor=south] at ({0.5*(\xsl+\xsr)},\ytop+0.55){$\phi_{\mathrm{s}}$ (sample)};
\draw (\xsr,\ytop) -- (\xrr,\ytop);
\node[railnode] at (\xsr+1.1,\ytop){};
\node[nlab, anchor=south, yshift=2pt] at (\xsr+1.1,\ytop){top plate \texttt{TP}};

% ===== one binary-weighted slice =====================================
% #1 x  #2 capacitor label  #3 bit index
\newcommand{\sarbit}[3]{%
  \draw (#1,\ytop) to[C, l_={#2}] (#1,\ybot);
  \draw (#1,\ybot) -- (#1,\ypiv);
  \node[railnode] at (#1,\ypiv){};
  \node[nlab, anchor=west, xshift=2pt] at (#1+0.06,\ypiv+0.34){\texttt{b}\sarbv{#3}};
  % lever, thrown toward the VREF contact
  \draw (#1,\ypiv) -- (#1-0.40,\ycon+0.10);
  % VREF contact and its drop to the reference rail
  \draw (#1-0.62,\ycon) -- (#1-0.28,\ycon);
  \draw (#1-0.45,\ycon) -- (#1-0.45,\yref);
  \node[railnode] at (#1-0.45,\yref){};
  % ground contact
  \draw (#1+0.28,\ycon) -- (#1+0.62,\ycon);
  \draw (#1+0.45,\ycon) -- (#1+0.45,\ycon-0.35) node[ground]{};
}
\sarbit{\xB}{$C$}{9}
\sarbit{\xC}{$C/2$}{8}
\sarbit{\xD}{$C/4$}{7}
\sarbit{\xF}{$C/512$}{0}

% ---- elided middle bits ---------------------------------------------
\node[font=\large] at (\xE,\ybot+0.9){$\cdots$};
\node[font=\large] at (\xE,\ycon+0.15){$\cdots$};
\node[note, anchor=north, align=center] at (\xE,\ybot+0.45)
  {bits 6--1:\\six identical\\slices, $C/8$--$C/256$};

% ---- terminating dummy, bottom plate hard-grounded ------------------
\draw (\xG,\ytop) to[C, l_={$C/512$}] (\xG,\ybot);
\draw (\xG,\ybot) -- (\xG,\ycon-0.35) node[ground]{};
\node[note, anchor=west, xshift=3pt] at (\xG+0.15,\ypiv+0.1)
  {terminating dummy\\(no bit switch)};

% ---- reference rail -------------------------------------------------
\draw[thick] (\xB-1.6,\yref) -- (\xF+0.2,\yref);
\node[nlab, anchor=east, xshift=-2pt] at (\xB-1.6,\yref){$V_{\mathrm{REF}}$};

% ===== comparator ====================================================
\node[op amp, noinv input up, anchor=+] (K) at (12.2,\ytop){};
\node[font=\scriptsize] at ($(K.center)+(-0.15,0)$) {CMP};
\draw (\xrr,\ytop) -- (K.+);
\draw (K.-) -- ++(-1.15,0) -- ++(0,-1.25) node[ground]{};
\node[nlab, anchor=east, xshift=-2pt] at ($(K.-)+(-1.15,-0.75)$){$V_{\mathrm{CM}}$};

% ===== SAR logic block ===============================================
\draw[thick, rounded corners=2pt] (0,-8.1) rectangle (9.75,-5.9);
\node[align=center] at (4.875,-7.0)
  {successive-approximation register\\and control logic};

% comparator decision into the logic
\draw[-{Stealth[length=4pt]}] (K.out) -- (14.7,0 |- K.out) -- (14.7,-6.35) -- (9.75,-6.35);
\node[nlab, anchor=south, yshift=1pt] at (12.2,-6.35){\texttt{cmp}};

% bit word back up to the switch bank
\draw[bus, -{Stealth[length=5pt]}] (4.875,-5.9) -- (4.875,-5.15);
\node[nlab, anchor=west, xshift=3pt] at (4.95,-5.5){\texttt{b}\sarbv{9:0}};

% clock and reset
\draw[-{Stealth[length=4pt]}] (-2.1,-6.6) -- (0,-6.6);
\node[port, anchor=east] at (-2.2,-6.6){Clock\_Input};
\draw[-{Stealth[length=4pt]}] (-2.1,-7.5) -- (0,-7.5);
\node[port, anchor=east] at (-2.2,-7.5){Reset\_Input};

% data and ready
\draw[bus, -{Stealth[length=5pt]}] (6.5,-8.1) -- (6.5,-9.15);
\draw (6.25,-8.75) -- (6.75,-8.35);
\node[nlab, anchor=east, xshift=-1pt] at (6.28,-8.5){10};
\node[port, anchor=north] at (6.5,-9.25){Data\_Output\sarbv{9:0}};
\draw[-{Stealth[length=4pt]}] (2.6,-8.1) -- (2.6,-9.15);
\node[port, anchor=north] at (2.6,-9.25){Ready\_Output};

\end{circuitikz}
  }
  \caption[{Generic charge-redistribution SAR ADC architecture.}]{Generic single-ended charge-redistribution SAR ADC, drawn to explain the architecture of the 10-bit converter rather than to reproduce its implemented transistor-level schematic. The input is sampled through $\phi_{\mathrm{s}}$ onto the top plate of a binary-weighted array; the logic then switches each bottom plate from $V_{\mathrm{REF}}$ to ground one bit at a time, most significant first, and keeps or discards the trial from the comparator decision, so the conversion takes ten comparator cycles and the top plate converges to $V_{\mathrm{CM}}$. Bits 6--1 are elided: six further slices identical to those drawn sit between $C/4$ and $C/512$, and the terminating dummy $C/512$ carries no bit switch, which makes the total array $2C$ and the LSB weight $C/512$. Only the external port names are those of the implemented block; the sampling network, comparator topology, clock generation and reset behaviour are representative, and no device dimensions are stated.}
  \label{fig:adc-sar}
\end{figure}
